\Chapter{79}

\Verse{1}
{
    \zA{话说宝玉才祭完了晴雯，只听花阴中有个人声，倒吓了一跳。}
}
{
    \eA{[English source not present in the checked-in Bencraft/Joly files.]}
}
{
}

\Verse{2}
{
    \zA{细看不是别人， 却是黛玉，满面含笑，口内说道：“好新奇的祭文!可与《曹娥碑》并传了。”}
}
{
    \eA{[English source not present in the checked-in Bencraft/Joly files.]}
}
{
}

\Verse{3}
{
    \zA{宝 玉听了，不觉红了脸，笑答道：“我想着世上这些祭文，都过于熟烂了，所以改个 新样。}
}
{
    \eA{[English source not present in the checked-in Bencraft/Joly files.]}
}
{
}

\Verse{4}
{
    \zA{原不过是我一时的玩意儿，谁知被你听见了。}
}
{
    \eA{[English source not present in the checked-in Bencraft/Joly files.]}
}
{
}

\Verse{5}
{
    \zA{有什么大使不得的，何不改削 改削？”}
}
{
    \eA{[English source not present in the checked-in Bencraft/Joly files.]}
}
{
}

\Verse{6}
{
    \zA{黛玉道：“原稿在那里？}
}
{
    \eA{[English source not present in the checked-in Bencraft/Joly files.]}
}
{
}

\Verse{7}
{
    \zA{倒要细细的看看。}
}
{
    \eA{[English source not present in the checked-in Bencraft/Joly files.]}
}
{
}

\Verse{8}
{
    \zA{长篇大论，不知说的是什么。}
}
{
    \eA{[English source not present in the checked-in Bencraft/Joly files.]}
}
{
}

\Verse{9}
{
    \zA{只听见中间两句，什么‘红绡帐里，公子情深；黄土陇中，女儿命薄’，这一联意 思却好。}
}
{
    \eA{[English source not present in the checked-in Bencraft/Joly files.]}
}
{
}

\Verse{10}
{
    \zA{只是‘红绡帐里’未免俗滥些。}
}
{
    \eA{[English source not present in the checked-in Bencraft/Joly files.]}
}
{
}

\Verse{11}
{
    \zA{放着现成的真事，为什么不用？”}
}
{
    \eA{[English source not present in the checked-in Bencraft/Joly files.]}
}
{
}

\Verse{12}
{
    \zA{宝玉忙 问：“什么现成的真事？”}
}
{
    \eA{[English source not present in the checked-in Bencraft/Joly files.]}
}
{
}

\Verse{13}
{
    \zA{黛玉笑道：“咱们如今都系霞彩纱糊的窗？}
}
{
    \eA{[English source not present in the checked-in Bencraft/Joly files.]}
}
{
}

\Verse{14}
{
    \zA{，何不说‘茜 纱窗下，公子多情’呢？”}
}
{
    \eA{[English source not present in the checked-in Bencraft/Joly files.]}
}
{
}

\Verse{15}
{
    \zA{宝玉听了，不禁跌脚笑道：“好极，好极!到底是你想 得出，说得出。}
}
{
    \eA{[English source not present in the checked-in Bencraft/Joly files.]}
}
{
}

\Verse{16}
{
    \zA{可知天下古今现成的好景好事尽多，只是我们愚人想不出来罢了。}
}
{
    \eA{[English source not present in the checked-in Bencraft/Joly files.]}
}
{
}

\Verse{17}
{
    \zA{但只一件：虽然这一改新妙之极，却是你在这里住着还可以，我实不敢当。”}
}
{
    \eA{[English source not present in the checked-in Bencraft/Joly files.]}
}
{
}

\Verse{18}
{
    \zA{说着， 又连说“不敢”。}
}
{
    \eA{[English source not present in the checked-in Bencraft/Joly files.]}
}
{
}

\Verse{19}
{
    \zA{黛玉笑道：“何妨？}
}
{
    \eA{[English source not present in the checked-in Bencraft/Joly files.]}
}
{
}

\Verse{20}
{
    \zA{我的窗即可为你之窗，何必如此分晰，也太 生疏了。}
}
{
    \eA{[English source not present in the checked-in Bencraft/Joly files.]}
}
{
}

\Verse{21}
{
    \zA{古人异姓陌路，尚然‘肥马轻裘，敝之无憾’，何况咱们？”}
}
{
    \eA{[English source not present in the checked-in Bencraft/Joly files.]}
}
{
}

\Verse{22}
{
    \zA{宝玉笑道： “论交道，不在‘肥马轻裘’，即黄金白璧亦不当锱铢较量。}
}
{
    \eA{[English source not present in the checked-in Bencraft/Joly files.]}
}
{
}

\Verse{23}
{
    \zA{倒是这唐突闺阁上头， 却万万使不得的。}
}
{
    \eA{[English source not present in the checked-in Bencraft/Joly files.]}
}
{
}

\Verse{24}
{
    \zA{如今我索性将‘公子’‘女儿’改去，竟算是你诔他的倒妙。}
}
{
    \eA{[English source not present in the checked-in Bencraft/Joly files.]}
}
{
}

\Verse{25}
{
    \zA{况 且素日你又待他甚厚，所以宁可弃了这一篇文，万不可弃这‘茜纱’新句。}
}
{
    \eA{[English source not present in the checked-in Bencraft/Joly files.]}
}
{
}

\Verse{26}
{
    \zA{莫若改 作‘茜纱窗下，小姐多情；黄土陇中，丫鬟薄命’。}
}
{
    \eA{[English source not present in the checked-in Bencraft/Joly files.]}
}
{
}

\Verse{27}
{
    \zA{如此一改，虽与我不涉，我也 惬怀。”}
}
{
    \eA{[English source not present in the checked-in Bencraft/Joly files.]}
}
{
}

\Verse{28}
{
    \zA{黛玉笑道：“他又不是我的丫头，何用此话？}
}
{
    \eA{[English source not present in the checked-in Bencraft/Joly files.]}
}
{
}

\Verse{29}
{
    \zA{况且‘小姐’‘丫鬟’，亦 不典雅。}
}
{
    \eA{[English source not present in the checked-in Bencraft/Joly files.]}
}
{
}

\Verse{30}
{
    \zA{等得紫鹃死了，我再如此说，还不算迟呢。”}
}
{
    \eA{[English source not present in the checked-in Bencraft/Joly files.]}
}
{
}

\Verse{31}
{
    \zA{宝玉听了笑道：“这是何苦， 又咒他。”}
}
{
    \eA{[English source not present in the checked-in Bencraft/Joly files.]}
}
{
}

\Verse{32}
{
    \zA{黛玉笑道：“是你要咒的，并不是我说的。”}
}
{
    \eA{[English source not present in the checked-in Bencraft/Joly files.]}
}
{
}

\Verse{33}
{
    \zA{宝玉说：“我又有了，这 一改恰就妥当了：莫若说‘茜纱窗下，我本无缘；黄土陇中，卿何薄命！’”}
}
{
    \eA{[English source not present in the checked-in Bencraft/Joly files.]}
}
{
}

\Verse{34}
{
    \zA{黛玉听了，陡然变色。}
}
{
    \eA{[English source not present in the checked-in Bencraft/Joly files.]}
}
{
}

\Verse{35}
{
    \zA{虽有无限狐疑，外面却不肯露出，反连忙含笑点头称妙， 说：“果然改得好。}
}
{
    \eA{[English source not present in the checked-in Bencraft/Joly files.]}
}
{
}

\Verse{36}
{
    \zA{再不必乱改了，快去干正经事罢。}
}
{
    \eA{[English source not present in the checked-in Bencraft/Joly files.]}
}
{
}

\Verse{37}
{
    \zA{刚才太太打发人叫你，说明 儿一早过大舅母那边去呢。}
}
{
    \eA{[English source not present in the checked-in Bencraft/Joly files.]}
}
{
}

\Verse{38}
{
    \zA{你二姐姐已有人家求准了，所以叫你们过去呢。”}
}
{
    \eA{[English source not present in the checked-in Bencraft/Joly files.]}
}
{
}

\Verse{39}
{
    \zA{宝玉 忙道：“何必如此忙？}
}
{
    \eA{[English source not present in the checked-in Bencraft/Joly files.]}
}
{
}

\Verse{40}
{
    \zA{我身上也不大好，明儿还未必能去呢。”}
}
{
    \eA{[English source not present in the checked-in Bencraft/Joly files.]}
}
{
}

\Verse{41}
{
    \zA{黛玉道：“又来了。}
}
{
    \eA{[English source not present in the checked-in Bencraft/Joly files.]}
}
{
}

\Verse{42}
{
    \zA{我劝你把脾气改改罢。}
}
{
    \eA{[English source not present in the checked-in Bencraft/Joly files.]}
}
{
}

\Verse{43}
{
    \zA{一年大，二年小，……”一面说话一面咳嗽起来。}
}
{
    \eA{[English source not present in the checked-in Bencraft/Joly files.]}
}
{
}

\Verse{44}
{
    \zA{宝玉忙道： “这里风冷，咱们只顾站着，凉着呢可不是玩的，快回去罢。”}
}
{
    \eA{[English source not present in the checked-in Bencraft/Joly files.]}
}
{
}

\Verse{45}
{
    \zA{黛玉道：“我也家 去歇息了，明儿再见罢。”}
}
{
    \eA{[English source not present in the checked-in Bencraft/Joly files.]}
}
{
}

\Verse{46}
{
    \zA{说着，便自取路去了。}
}
{
    \eA{[English source not present in the checked-in Bencraft/Joly files.]}
}
{
}

\Verse{47}
{
    \zA{宝玉只得闷闷的转步，忽想起黛 玉无人随伴，忙命小丫头子跟送回去。}
}
{
    \eA{[English source not present in the checked-in Bencraft/Joly files.]}
}
{
}

\Verse{48}
{
    \zA{自己到了怡红院中，果有王夫人打发嬷嬷们 来，吩咐他明日一早过贾赦这边来，与方才黛玉之言相对。}
}
{
    \eA{[English source not present in the checked-in Bencraft/Joly files.]}
}
{
}

\Verse{49}
{
    \zA{原来贾赦已将迎春许与孙家了。}
}
{
    \eA{[English source not present in the checked-in Bencraft/Joly files.]}
}
{
}

\Verse{50}
{
    \zA{这孙家乃是大同府人氏，祖上系军官出身，乃 当日宁荣府中之门生，算来亦系至交。}
}
{
    \eA{[English source not present in the checked-in Bencraft/Joly files.]}
}
{
}

\Verse{51}
{
    \zA{如今孙家只有一人在京，现袭指挥之职。}
}
{
    \eA{[English source not present in the checked-in Bencraft/Joly files.]}
}
{
}

\Verse{52}
{
    \zA{此 人名唤孙绍祖，生得相貌魁梧，体格健壮，弓马娴熟，应酬权变，年纪未满三十， 且又家资饶富，现在兵部候缺题升。}
}
{
    \eA{[English source not present in the checked-in Bencraft/Joly files.]}
}
{
}

\Verse{53}
{
    \zA{因未曾娶妻，贾赦见是世交子侄，且人品家当 都相称合，遂择为东床娇婿。}
}
{
    \eA{[English source not present in the checked-in Bencraft/Joly files.]}
}
{
}

\Verse{54}
{
    \zA{亦曾回明贾母，贾母心中却不大愿意，但想儿女之事， 自有天意，况且他亲父主张，何必出头多事？}
}
{
    \eA{[English source not present in the checked-in Bencraft/Joly files.]}
}
{
}

\Verse{55}
{
    \zA{因此只说“知道了”三字，馀不多及。}
}
{
    \eA{[English source not present in the checked-in Bencraft/Joly files.]}
}
{
}

\Verse{56}
{
    \zA{贾政又深恶孙家，虽是世交，不过是他祖父当日希慕宁荣之势，有不能了结之事挽 拜在门下的，并非诗礼名族之裔。}
}
{
    \eA{[English source not present in the checked-in Bencraft/Joly files.]}
}
{
}

\Verse{57}
{
    \zA{因此，倒劝谏过两次，无奈贾赦不听，也只得罢 了。}
}
{
    \eA{[English source not present in the checked-in Bencraft/Joly files.]}
}
{
}

\Verse{58}
{
    \zA{宝玉却未曾会过这孙绍祖一面的，次日只得过去，聊以塞责。}
}
{
    \eA{[English source not present in the checked-in Bencraft/Joly files.]}
}
{
}

\Verse{59}
{
    \zA{只听见那娶亲的 日子甚近，不过今年就要过门的，又见邢夫人等回了贾母，将迎春接出大观园去， 越发扫兴。}
}
{
    \eA{[English source not present in the checked-in Bencraft/Joly files.]}
}
{
}

\Verse{60}
{
    \zA{每每痴痴呆呆的，不知作何消遣。}
}
{
    \eA{[English source not present in the checked-in Bencraft/Joly files.]}
}
{
}

\Verse{61}
{
    \zA{又听说要陪四个丫头过去，更又跌足 道：“从今后这世上又少了五个清净人了！”}
}
{
    \eA{[English source not present in the checked-in Bencraft/Joly files.]}
}
{
}

\Verse{62}
{
    \zA{因此天天到紫菱洲一带地方徘徊瞻顾。}
}
{
    \eA{[English source not present in the checked-in Bencraft/Joly files.]}
}
{
}

\Verse{63}
{
    \zA{见其轩窗寂寞，屏帐？}
}
{
    \eA{[English source not present in the checked-in Bencraft/Joly files.]}
}
{
}

\Verse{64}
{
    \zA{然，不过只有几个该班上夜的老妪。}
}
{
    \eA{[English source not present in the checked-in Bencraft/Joly files.]}
}
{
}

\Verse{65}
{
    \zA{再看那岸上的蓼花苇叶， 也都觉摇摇落落，似有追忆故人之态，迥非素常逞妍斗色可比。}
}
{
    \eA{[English source not present in the checked-in Bencraft/Joly files.]}
}
{
}

\Verse{66}
{
    \zA{所以情不自禁，乃 信口吟成一歌曰： 池塘一夜秋风冷，吹散芰荷红玉影。}
}
{
    \eA{[English source not present in the checked-in Bencraft/Joly files.]}
}
{
}

\Verse{67}
{
    \zA{蓼花菱叶不胜悲，重露繁霜压纤梗。}
}
{
    \eA{[English source not present in the checked-in Bencraft/Joly files.]}
}
{
}

\Verse{68}
{
    \zA{不闻永昼敲棋声，燕泥点点污棋枰。}
}
{
    \eA{[English source not present in the checked-in Bencraft/Joly files.]}
}
{
}

\Verse{69}
{
    \zA{古人惜别怜朋友，况我今当手足情! 宝玉方才吟罢，忽闻背后有人笑道：“你又发什么呆呢？”}
}
{
    \eA{[English source not present in the checked-in Bencraft/Joly files.]}
}
{
}

\Verse{70}
{
    \zA{宝玉回头忙看是谁， 原来是香菱。}
}
{
    \eA{[English source not present in the checked-in Bencraft/Joly files.]}
}
{
}

\Verse{71}
{
    \zA{宝玉忙转身笑问道：“我的姐姐，你这会子跑到这里来做什么？}
}
{
    \eA{[English source not present in the checked-in Bencraft/Joly files.]}
}
{
}

\Verse{72}
{
    \zA{许多 日子也不进来逛逛。”}
}
{
    \eA{[English source not present in the checked-in Bencraft/Joly files.]}
}
{
}

\Verse{73}
{
    \zA{香菱拍手笑嘻嘻的说道：“我何曾不要来。}
}
{
    \eA{[English source not present in the checked-in Bencraft/Joly files.]}
}
{
}

\Verse{74}
{
    \zA{如今你哥哥回来 了，那里比先时自由自在的了？}
}
{
    \eA{[English source not present in the checked-in Bencraft/Joly files.]}
}
{
}

\Verse{75}
{
    \zA{才刚我们太太使人找你凤姐姐去，竟没有找着，说 往园子里来了。}
}
{
    \eA{[English source not present in the checked-in Bencraft/Joly files.]}
}
{
}

\Verse{76}
{
    \zA{我听见这个话，我就讨了这个差进来找他。}
}
{
    \eA{[English source not present in the checked-in Bencraft/Joly files.]}
}
{
}

\Verse{77}
{
    \zA{遇见他的丫头，说在稻 香村呢。}
}
{
    \eA{[English source not present in the checked-in Bencraft/Joly files.]}
}
{
}

\Verse{78}
{
    \zA{如今我往稻香村去，谁知又遇见了你。}
}
{
    \eA{[English source not present in the checked-in Bencraft/Joly files.]}
}
{
}

\Verse{79}
{
    \zA{我还要问你：袭人姐姐这几日可好？}
}
{
    \eA{[English source not present in the checked-in Bencraft/Joly files.]}
}
{
}

\Verse{80}
{
    \zA{怎么忽然把个晴雯姐姐也没了？}
}
{
    \eA{[English source not present in the checked-in Bencraft/Joly files.]}
}
{
}

\Verse{81}
{
    \zA{到底是什么病？}
}
{
    \eA{[English source not present in the checked-in Bencraft/Joly files.]}
}
{
}

\Verse{82}
{
    \zA{二姑娘搬出去的好快!你瞧瞧，这地 方一时间就空落落的了。”}
}
{
    \eA{[English source not present in the checked-in Bencraft/Joly files.]}
}
{
}

\Verse{83}
{
    \zA{宝玉只有一味答应，又让他同到怡红院去吃茶。}
}
{
    \eA{[English source not present in the checked-in Bencraft/Joly files.]}
}
{
}

\Verse{84}
{
    \zA{香菱道： “此刻竟不能，等找着琏二奶奶，说完了正经话再来。”}
}
{
    \eA{[English source not present in the checked-in Bencraft/Joly files.]}
}
{
}

\Verse{85}
{
    \zA{宝玉道：“什么正经话， 这般忙？”}
}
{
    \eA{[English source not present in the checked-in Bencraft/Joly files.]}
}
{
}

\Verse{86}
{
    \zA{香菱道：“为你哥哥娶嫂子的话，所以要紧。”}
}
{
    \eA{[English source not present in the checked-in Bencraft/Joly files.]}
}
{
}

\Verse{87}
{
    \zA{宝玉道：“正是说的是 那一家的好？}
}
{
    \eA{[English source not present in the checked-in Bencraft/Joly files.]}
}
{
}

\Verse{88}
{
    \zA{只听见吵嚷了这半年，今儿又说张家的好，明儿又要李家的，后儿又 议论王家的好。}
}
{
    \eA{[English source not present in the checked-in Bencraft/Joly files.]}
}
{
}

\Verse{89}
{
    \zA{这些人家的女儿，他也不知造了什么罪，叫人家好端端的议论。”}
}
{
    \eA{[English source not present in the checked-in Bencraft/Joly files.]}
}
{
}

\Verse{90}
{
    \zA{香菱道：“如今定了，可以不用拉扯别人家了。”}
}
{
    \eA{[English source not present in the checked-in Bencraft/Joly files.]}
}
{
}

\Verse{91}
{
    \zA{宝玉问道：“定了谁家的？”}
}
{
    \eA{[English source not present in the checked-in Bencraft/Joly files.]}
}
{
}

\Verse{92}
{
    \zA{香 菱道：“因你哥哥上次出门时，顺路到了个亲戚家去。}
}
{
    \eA{[English source not present in the checked-in Bencraft/Joly files.]}
}
{
}

\Verse{93}
{
    \zA{这门亲原是老亲，且又和我 们是同在户部挂名行商，也是数一数二的大门户。}
}
{
    \eA{[English source not present in the checked-in Bencraft/Joly files.]}
}
{
}

\Verse{94}
{
    \zA{前日说起来时，你们两府都也知 道的：合京城里，上至王侯，下至买卖人，都称他家是‘桂花夏家’。”}
}
{
    \eA{[English source not present in the checked-in Bencraft/Joly files.]}
}
{
}

\Verse{95}
{
    \zA{宝玉忙笑 道：“如何又称为‘桂花夏家’？”}
}
{
    \eA{[English source not present in the checked-in Bencraft/Joly files.]}
}
{
}

\Verse{96}
{
    \zA{香菱道：“本姓夏，非常的富贵。}
}
{
    \eA{[English source not present in the checked-in Bencraft/Joly files.]}
}
{
}

\Verse{97}
{
    \zA{其馀田地不 用说，单有几十顷地种着桂花，凡这长安那城里城外桂花局，俱是他家的，连宫里 一应陈设盆景，亦是他家供奉。}
}
{
    \eA{[English source not present in the checked-in Bencraft/Joly files.]}
}
{
}

\Verse{98}
{
    \zA{因此才有这个混号。}
}
{
    \eA{[English source not present in the checked-in Bencraft/Joly files.]}
}
{
}

\Verse{99}
{
    \zA{如今太爷也没了，只有老奶奶 带着一个亲生的姑娘过活，也并没有哥儿弟兄。}
}
{
    \eA{[English source not present in the checked-in Bencraft/Joly files.]}
}
{
}

\Verse{100}
{
    \zA{可惜他竟一门尽绝了后。”}
}
{
    \eA{[English source not present in the checked-in Bencraft/Joly files.]}
}
{
}

\Verse{101}
{
    \zA{宝玉忙 道：“咱们也别管他绝后不绝后，只是这姑娘可好？}
}
{
    \eA{[English source not present in the checked-in Bencraft/Joly files.]}
}
{
}

\Verse{102}
{
    \zA{你们大爷怎么就中意了？”}
}
{
    \eA{[English source not present in the checked-in Bencraft/Joly files.]}
}
{
}

\Verse{103}
{
    \zA{香 菱笑道：“一则是天缘，二来是‘情人眼里出西施’。}
}
{
    \eA{[English source not present in the checked-in Bencraft/Joly files.]}
}
{
}

\Verse{104}
{
    \zA{当年时又通家来往，从小儿 都在一处玩过。}
}
{
    \eA{[English source not present in the checked-in Bencraft/Joly files.]}
}
{
}

\Verse{105}
{
    \zA{叙亲是姑舅兄妹，又没嫌疑。}
}
{
    \eA{[English source not present in the checked-in Bencraft/Joly files.]}
}
{
}

\Verse{106}
{
    \zA{虽离了这几年，前儿一到他家，夏奶 奶又是没儿子的，一见了你哥哥出落的这么，又是哭，又是笑，竟比见了儿子的还 胜。}
}
{
    \eA{[English source not present in the checked-in Bencraft/Joly files.]}
}
{
}

\Verse{107}
{
    \zA{又令他兄妹相见。}
}
{
    \eA{[English source not present in the checked-in Bencraft/Joly files.]}
}
{
}

\Verse{108}
{
    \zA{谁知这姑娘出落的花朵似的了，在家里也读书写字，所以你 哥哥当时就一心看准了。}
}
{
    \eA{[English source not present in the checked-in Bencraft/Joly files.]}
}
{
}

\Verse{109}
{
    \zA{连当铺里老伙计们一群人，遭扰了人家三四日。}
}
{
    \eA{[English source not present in the checked-in Bencraft/Joly files.]}
}
{
}

\Verse{110}
{
    \zA{他们还留 多住几天，好容易苦辞，才放回家。}
}
{
    \eA{[English source not present in the checked-in Bencraft/Joly files.]}
}
{
}

\Verse{111}
{
    \zA{你哥哥一进门，就咕咕唧唧求我们太太去求亲。}
}
{
    \eA{[English source not present in the checked-in Bencraft/Joly files.]}
}
{
}

\Verse{112}
{
    \zA{我们太太原是见过的，又且门当户对，也依了。}
}
{
    \eA{[English source not present in the checked-in Bencraft/Joly files.]}
}
{
}

\Verse{113}
{
    \zA{和这里姨太太凤姑娘商议了打发人 去一说，就成了。}
}
{
    \eA{[English source not present in the checked-in Bencraft/Joly files.]}
}
{
}

\Verse{114}
{
    \zA{只是娶的日子太急，所以我们忙乱的很。}
}
{
    \eA{[English source not present in the checked-in Bencraft/Joly files.]}
}
{
}

\Verse{115}
{
    \zA{我也巴不得早些过来， 又添了一个做诗的人了。”}
}
{
    \eA{[English source not present in the checked-in Bencraft/Joly files.]}
}
{
}

\Verse{116}
{
    \zA{宝玉冷笑道：“虽如此说，但只我倒替你担心虑后呢。”}
}
{
    \eA{[English source not present in the checked-in Bencraft/Joly files.]}
}
{
}

\Verse{117}
{
    \zA{香菱道：“这是什么话？}
}
{
    \eA{[English source not present in the checked-in Bencraft/Joly files.]}
}
{
}

\Verse{118}
{
    \zA{我倒不懂了。”}
}
{
    \eA{[English source not present in the checked-in Bencraft/Joly files.]}
}
{
}

\Verse{119}
{
    \zA{宝玉笑道：“这有什么不懂的？}
}
{
    \eA{[English source not present in the checked-in Bencraft/Joly files.]}
}
{
}

\Verse{120}
{
    \zA{只怕再有个 人来，薛大哥就不肯疼你了。”}
}
{
    \eA{[English source not present in the checked-in Bencraft/Joly files.]}
}
{
}

\Verse{121}
{
    \zA{香菱听了，不觉红了脸，正色道：“这是怎么说？}
}
{
    \eA{[English source not present in the checked-in Bencraft/Joly files.]}
}
{
}

\Verse{122}
{
    \zA{素日咱们都是厮抬厮敬，今日忽然提起这些事来。}
}
{
    \eA{[English source not present in the checked-in Bencraft/Joly files.]}
}
{
}

\Verse{123}
{
    \zA{怪不得人人都说你是个亲近不得 的人。”}
}
{
    \eA{[English source not present in the checked-in Bencraft/Joly files.]}
}
{
}

\Verse{124}
{
    \zA{一面说，一面转身走了。}
}
{
    \eA{[English source not present in the checked-in Bencraft/Joly files.]}
}
{
}

\Verse{125}
{
    \zA{宝玉见他这样，便怅然如有所失，呆呆的站了半日，只得没精打彩，还入怡红 院来。}
}
{
    \eA{[English source not present in the checked-in Bencraft/Joly files.]}
}
{
}

\Verse{126}
{
    \zA{一夜不曾安睡，种种不宁。}
}
{
    \eA{[English source not present in the checked-in Bencraft/Joly files.]}
}
{
}

\Verse{127}
{
    \zA{次日便懒进饮食，身体发热。}
}
{
    \eA{[English source not present in the checked-in Bencraft/Joly files.]}
}
{
}

\Verse{128}
{
    \zA{也因近日抄检大观 园、逐司棋、别迎春、悲晴雯等羞辱、惊恐、悲凄所致，兼以风寒外感，遂致成疾， 卧床不起。}
}
{
    \eA{[English source not present in the checked-in Bencraft/Joly files.]}
}
{
}

\Verse{129}
{
    \zA{贾母听得如此，天天亲来看视。}
}
{
    \eA{[English source not present in the checked-in Bencraft/Joly files.]}
}
{
}

\Verse{130}
{
    \zA{王夫人心中自悔，不合因晴雯过于逼责 了他。}
}
{
    \eA{[English source not present in the checked-in Bencraft/Joly files.]}
}
{
}

\Verse{131}
{
    \zA{心中虽如此，脸上却不露出，只吩咐众奶娘等好生伏侍看守。}
}
{
    \eA{[English source not present in the checked-in Bencraft/Joly files.]}
}
{
}

\Verse{132}
{
    \zA{一日两次带进 医生来诊脉下药。}
}
{
    \eA{[English source not present in the checked-in Bencraft/Joly files.]}
}
{
}

\Verse{133}
{
    \zA{一月之后，方才渐渐的痊愈。}
}
{
    \eA{[English source not present in the checked-in Bencraft/Joly files.]}
}
{
}

\Verse{134}
{
    \zA{好生保养过百日，方许动荤腥油面， 方可出门行走。}
}
{
    \eA{[English source not present in the checked-in Bencraft/Joly files.]}
}
{
}

\Verse{135}
{
    \zA{这百日内，院门前皆不许到，只在屋里玩笑。}
}
{
    \eA{[English source not present in the checked-in Bencraft/Joly files.]}
}
{
}

\Verse{136}
{
    \zA{四五十天后，就把他 拘的火星乱迸，那里忍耐的住？}
}
{
    \eA{[English source not present in the checked-in Bencraft/Joly files.]}
}
{
}

\Verse{137}
{
    \zA{虽百般设法，无奈贾母王夫人执意不从，也只得罢 了。}
}
{
    \eA{[English source not present in the checked-in Bencraft/Joly files.]}
}
{
}

\Verse{138}
{
    \zA{因此，和些丫鬟们无所不至，恣意耍笑。}
}
{
    \eA{[English source not present in the checked-in Bencraft/Joly files.]}
}
{
}

\Verse{139}
{
    \zA{又听得薛蟠那里摆酒唱戏，热闹非常， 已娶亲入门。}
}
{
    \eA{[English source not present in the checked-in Bencraft/Joly files.]}
}
{
}

\Verse{140}
{
    \zA{闻得这夏家小姐十分俊俏，也略通文翰，宝玉恨不得就过去一见才好。}
}
{
    \eA{[English source not present in the checked-in Bencraft/Joly files.]}
}
{
}

\Verse{141}
{
    \zA{再过些时，又闻得迎春出了阁。}
}
{
    \eA{[English source not present in the checked-in Bencraft/Joly files.]}
}
{
}

\Verse{142}
{
    \zA{宝玉思及当时姊妹耳鬓厮磨，从今一别，纵得相逢， 必不得似先前这等亲热了。}
}
{
    \eA{[English source not present in the checked-in Bencraft/Joly files.]}
}
{
}

\Verse{143}
{
    \zA{眼前又不能去一望，真令人凄惶不尽。}
}
{
    \eA{[English source not present in the checked-in Bencraft/Joly files.]}
}
{
}

\Verse{144}
{
    \zA{少不得潜心忍耐， 暂同这些丫鬟们厮闹释闷，幸免贾政责备逼迫读书之难。}
}
{
    \eA{[English source not present in the checked-in Bencraft/Joly files.]}
}
{
}

\Verse{145}
{
    \zA{这百日内，只不曾拆毁了 怡红院，和这些丫头们无法无天，凡世上所无之事，都玩耍出来，如今且不消细说。}
}
{
    \eA{[English source not present in the checked-in Bencraft/Joly files.]}
}
{
}

\Verse{146}
{
    \zA{且说香菱自那日抢白了宝玉之后，自为宝玉有意唐突，“从此倒要远避他些才 好。”}
}
{
    \eA{[English source not present in the checked-in Bencraft/Joly files.]}
}
{
}

\Verse{147}
{
    \zA{因此，以后连大观园也不轻易进来了。}
}
{
    \eA{[English source not present in the checked-in Bencraft/Joly files.]}
}
{
}

\Verse{148}
{
    \zA{日日忙乱着薛蟠娶过亲，因为得了护 身符，自己身上分去责任，到底比这样安静些；二则又知是个有才有貌的佳人，自
        然是典雅和平的：因此，心里盼过门的日子比薛蟠还急十倍呢。}
}
{
    \eA{[English source not present in the checked-in Bencraft/Joly files.]}
}
{
}

\Verse{149}
{
    \zA{好容易盼得一日娶 过来，他便十分殷勤小心伏侍。}
}
{
    \eA{[English source not present in the checked-in Bencraft/Joly files.]}
}
{
}

\Verse{150}
{
    \zA{原来这夏家小姐今年方十七岁，生得亦颇有姿色，亦颇识得几个字。}
}
{
    \eA{[English source not present in the checked-in Bencraft/Joly files.]}
}
{
}

\Verse{151}
{
    \zA{若论心里 的丘壑泾渭，颇步熙凤的后尘。}
}
{
    \eA{[English source not present in the checked-in Bencraft/Joly files.]}
}
{
}

\Verse{152}
{
    \zA{只吃亏了一件：从小时父亲去世的早，又无同胞兄 弟，寡母独守此女，娇养溺爱，不啻珍宝，凡女儿一举一动，他母亲皆百依百顺，
        因此未免酿成个盗跖的情性：自己尊若菩萨，他人秽如粪土；外具花柳之姿，内秉 风雷之性。}
}
{
    \eA{[English source not present in the checked-in Bencraft/Joly files.]}
}
{
}

\Verse{153}
{
    \zA{在家里和丫鬟们使性赌气、轻骂重打的。}
}
{
    \eA{[English source not present in the checked-in Bencraft/Joly files.]}
}
{
}

\Verse{154}
{
    \zA{今儿出了阁，自为要作当家的 奶奶，比不得做女儿时腼腆温柔，须要拿出威风来才钤压得住人。}
}
{
    \eA{[English source not present in the checked-in Bencraft/Joly files.]}
}
{
}

\Verse{155}
{
    \zA{况且见薛蟠气质 刚硬，举止骄奢，若不趁热灶一气炮制，将来必不能自竖旗帜矣。}
}
{
    \eA{[English source not present in the checked-in Bencraft/Joly files.]}
}
{
}

\Verse{156}
{
    \zA{又见有香菱这等 一个才貌俱全的爱妾在室，越发添了“宋太祖灭南唐”之意。}
}
{
    \eA{[English source not present in the checked-in Bencraft/Joly files.]}
}
{
}

\Verse{157}
{
    \zA{因他家多桂花，他小 名就叫做金桂。}
}
{
    \eA{[English source not present in the checked-in Bencraft/Joly files.]}
}
{
}

\Verse{158}
{
    \zA{他在家时，不许人口中带出“金”“桂”二字来，凡有不留心误道 一字者，他便定要苦打重罚才罢。}
}
{
    \eA{[English source not present in the checked-in Bencraft/Joly files.]}
}
{
}

\Verse{159}
{
    \zA{他因想“桂花”二字是禁止不住的，须得另换一 名，想桂花曾有广寒嫦娥之说，便将桂花改为“嫦娥花”，又寓自己身分。}
}
{
    \eA{[English source not present in the checked-in Bencraft/Joly files.]}
}
{
}

\Verse{160}
{
    \zA{如今薛 蟠本是个怜新弃旧的人，且是有酒胆、无饭力的，如今得了这一个妻子，正在新鲜 兴头上，凡事未免尽让他些。}
}
{
    \eA{[English source not present in the checked-in Bencraft/Joly files.]}
}
{
}

\Verse{161}
{
    \zA{那夏金桂见是这般形景，便也试着一步紧似一步。}
}
{
    \eA{[English source not present in the checked-in Bencraft/Joly files.]}
}
{
}

\Verse{162}
{
    \zA{一 月之中，二人气概都还相平；至两月之后，便觉薛蟠的气概渐次的低矮了下去。}
}
{
    \eA{[English source not present in the checked-in Bencraft/Joly files.]}
}
{
}

\Verse{163}
{
    \zA{一日，薛蟠酒后，不知要行何事，先和金桂商议。}
}
{
    \eA{[English source not present in the checked-in Bencraft/Joly files.]}
}
{
}

\Verse{164}
{
    \zA{金桂执意不从，薛蟠便忍不 住，便发了几句话，赌气自行了。}
}
{
    \eA{[English source not present in the checked-in Bencraft/Joly files.]}
}
{
}

\Verse{165}
{
    \zA{金桂便哭的如醉人一般，茶汤不进，装起病来， 请医疗治。}
}
{
    \eA{[English source not present in the checked-in Bencraft/Joly files.]}
}
{
}

\Verse{166}
{
    \zA{医生又说：“气血相逆，当进宽胸顺气之剂。”}
}
{
    \eA{[English source not present in the checked-in Bencraft/Joly files.]}
}
{
}

\Verse{167}
{
    \zA{薛姨妈恨得骂了薛蟠一 顿，说：“如今娶了亲，眼前抱儿子了，还是这么胡闹!人家凤凰似的，好容易养
        了一个女儿，比花朵儿还轻巧，原看的你是个人物，才给你做媳妇。}
}
{
    \eA{[English source not present in the checked-in Bencraft/Joly files.]}
}
{
}

\Verse{168}
{
    \zA{你不说收了心， 安分守己，一心一计，和和气气的过日子，还是这么胡闹，喝了黄汤折磨人家。}
}
{
    \eA{[English source not present in the checked-in Bencraft/Joly files.]}
}
{
}

\Verse{169}
{
    \zA{这 会子花钱吃药白遭心。”}
}
{
    \eA{[English source not present in the checked-in Bencraft/Joly files.]}
}
{
}

\Verse{170}
{
    \zA{一席话说的薛蟠后悔不迭，反来安慰金桂。}
}
{
    \eA{[English source not present in the checked-in Bencraft/Joly files.]}
}
{
}

\Verse{171}
{
    \zA{金桂见婆婆如 此说，越发得了意，更装出些张致来，不理薛蟠。}
}
{
    \eA{[English source not present in the checked-in Bencraft/Joly files.]}
}
{
}

\Verse{172}
{
    \zA{薛蟠没了主意，惟有自软而已。}
}
{
    \eA{[English source not present in the checked-in Bencraft/Joly files.]}
}
{
}

\Verse{173}
{
    \zA{好容易十天半月之后，才渐渐的哄转过金桂的心来。}
}
{
    \eA{[English source not present in the checked-in Bencraft/Joly files.]}
}
{
}

\Verse{174}
{
    \zA{自此，便加一倍小心，气概不免又矮了半截下来。}
}
{
    \eA{[English source not present in the checked-in Bencraft/Joly files.]}
}
{
}

\Verse{175}
{
    \zA{那金桂见丈夫旗纛渐倒，婆 婆良善，也就渐渐的持戈试马。}
}
{
    \eA{[English source not present in the checked-in Bencraft/Joly files.]}
}
{
}

\Verse{176}
{
    \zA{先时不过挟制薛蟠；后来倚娇作媚，将及薛姨妈； 后将至宝钗。}
}
{
    \eA{[English source not present in the checked-in Bencraft/Joly files.]}
}
{
}

\Verse{177}
{
    \zA{宝钗久察其不轨之心，每每随机应变，暗以言语弹压其志。}
}
{
    \eA{[English source not present in the checked-in Bencraft/Joly files.]}
}
{
}

\Verse{178}
{
    \zA{金桂知其 不可犯，便欲寻隙，苦得无隙可乘，倒只好曲意俯就。}
}
{
    \eA{[English source not present in the checked-in Bencraft/Joly files.]}
}
{
}

\Verse{179}
{
    \zA{一日，金桂无事，因和香菱 闲谈，问香菱家乡父母。}
}
{
    \eA{[English source not present in the checked-in Bencraft/Joly files.]}
}
{
}

\Verse{180}
{
    \zA{香菱皆答“忘记”，金桂便不悦，说有意欺瞒了他。}
}
{
    \eA{[English source not present in the checked-in Bencraft/Joly files.]}
}
{
}

\Verse{181}
{
    \zA{因问： “‘香菱’二字是谁起的？”}
}
{
    \eA{[English source not present in the checked-in Bencraft/Joly files.]}
}
{
}

\Verse{182}
{
    \zA{香菱便答道：“姑娘起的。”}
}
{
    \eA{[English source not present in the checked-in Bencraft/Joly files.]}
}
{
}

\Verse{183}
{
    \zA{金桂冷笑道：“人人都 说姑娘通，只这一个名字就不通。”}
}
{
    \eA{[English source not present in the checked-in Bencraft/Joly files.]}
}
{
}

\Verse{184}
{
    \zA{香菱忙笑道：“奶奶若说姑娘不通，奶奶没合 姑娘讲究过。}
}
{
    \eA{[English source not present in the checked-in Bencraft/Joly files.]}
}
{
}

\Verse{185}
{
    \zA{说起来，他的学问，连咱们姨老爷常时还夸的呢。”}
}
{
    \eA{[English source not present in the checked-in Bencraft/Joly files.]}
}
{
}

\Verse{186}
{
    \zA{欲知香菱说出何话，且听下回分解。}
}
{
    \eA{[English source not present in the checked-in Bencraft/Joly files.]}
}
{
}

\Verse{187}
{
    \zA{(本章完)}
}
{
    \eA{[English source not present in the checked-in Bencraft/Joly files.]}
}
{
}
